% ===================================================================
%  TAFEL Nr. 10 — Das Meer. --- Der Hafen.
%  Traduction 2026 d'apres texto/io, controlee sur texto/fr et
%  texto/en. Consigne : voir texto/de/10-tafel-01.tex.
%
%  Le francais decale sa numerotation d'un cran a partir de l'alinea 8
%  (il ecrit « 9. --- » sur la cle t10-08-1). C'est une coquille de
%  l'imprimeur francais, conservee dans la transcription ; l'allemand
%  suit l'ido, qui compte juste.
%
%  LE VOCABULAIRE DE LA MER EST BAS-ALLEMAND, ET C'EST POUR CELA
%  QU'IL EST ANGLAIS AUSSI. Boot, Mast, Kiel, Deck, Anker, Tau, Ebbe,
%  Flut, Werft, Dock, Bord, Möwe, Schiff : ce sont les mots de la
%  Hanse, montes du nord dans l'allemand ecrit, et ce sont mot pour
%  mot ceux de l'anglais — boat, mast, keel, deck, anchor, dock,
%  board, ship. Les deux langues ne se les sont pas prêtés : elles les
%  ont pris au meme commerce, sur la meme mer, en meme temps.
%
%  MAIS LA FLOTTE DE GUERRE DE CETTE PAGE EST DE 1895, ET ELLE A ETE
%  NOMMEE EXPRES. Quand l'Empire s'est donne une marine, l'etat-major
%  a refuse les termes internationaux et fait des composes, un par
%  type de navire :
%     Schlachtschiff      le cuirasse         (9)
%     Panzerkreuzer       le croiseur cuirasse (6)
%     Torpedoboot         le torpilleur       (8)
%     Unterseeboot, U-Boot le sous-marin      (10)
%     Kanonenboot         la canonniere       (7)
%     Truppentransporter  le transport        (12)
%     Schwimmdock         le dock flottant    (27)
%     Leuchtturm          le phare            (93)
%  C'est la quatrieme campagne de nomination que la serie rencontre,
%  apres Jahn (1811), le Sprachverein (1890) et les deux dernieres
%  colonnes de la premiere serie. Chaque navire de la planche porte un
%  compose allemand.
%
%  ET CHAQUE HOMME QUI EST DESSUS PORTE UN TITRE LATIN : Admiral,
%  Kommandant, Leutnant, Kapitän, Offizier, Quartiermeister, Matrose,
%  Lotse, Pilot. Pas un n'a ete traduit. LE NAVIRE EST ALLEMAND,
%  L'HOMME QUI LE COMMANDE EST LATIN — c'est la regle du tableau 1
%  appliquee a une marine : l'objet est indigene, le grade est
%  importe, et le grade est importe parce que la hierarchie l'est.
%
%  ENFIN LE « U-BOOT » REPOND A LA QUESTION QUE LE TABLEAU 9 AVAIT
%  LAISSEE OUVERTE. Le tableau 2 disait qu'un mot decrete ne dure que
%  le temps de celui qui decrete ; le tableau 9 qu'un mot voyage quand
%  la chose voyage. « U-Boot » est un mot DECRETE qui a VOYAGE : il
%  est reparti en anglais tel quel, « U-boat », vingt ans apres avoir
%  ete fabrique dans un bureau de Berlin. Les deux regles ne se
%  contredisent pas, elles se completent :
%
%     UN MOT DECRETE VOYAGE COMME UN AUTRE, A CONDITION QUE LA CHOSE
%     QU'IL NOMME VOYAGE AVEC LUI. Ce qui ne voyage pas, ce n'est pas
%     le mot fabrique, c'est le mot fabrique POUR UNE CHOSE QUI ETAIT
%     DEJA LA.
%
%  Le calendrier du tableau 3 avait echoue parce que les mois etaient
%  la ; « Turnen » a tenu chez soi et n'est parti nulle part parce que
%  la gymnastique etait partout ; « Gletscher » et « U-Boot » sont
%  partis parce que le glacier et le sous-marin arrivaient.
% ===================================================================

%%K t10-apar-1 apar
\VUsaut{4.0mm}
\VUtitre{15.0pt}{60}{TAFEL Nr. 10}
\VUsaut{3.0mm}
\VUpk{11.4pt}{40}{Das Meer. --- Der Hafen.}
\VUsaut{3.0mm}

%%K t10-01-1 p
1. --- Im Sommer suchen die einen Ruhe und kühle Luft im Gebirge, die anderen fahren
lieber an die Küste. Das haben wir letztes Jahr gemacht, weil wir den \VUgras{Ozean}
\textsuperscript{(1)} sehr lieben.

%%K t10-02-1 p
2. --- Ich schaue für mein Leben gern auf diese riesige Wasserfläche, die sich bis
zum \VUgras{Horizont} \textsuperscript{(2)} erstreckt, und sehe zu, wie die großen
\VUgras{Dampfer} \textsuperscript{(3)} zu einer langen \VUgras{Überfahrt} auslaufen,
mit den Reisenden nach Amerika an Bord.

%%K t10-03-1 p
3. --- Mein Vater, der weiß, was mir gefällt, sucht deshalb für die Ferien immer
einen sehr ruhigen Strand in der Nähe eines unserer großen \VUgras{Kriegshäfen} aus.

%%K t10-03-2 p
Bei unserem letzten Aufenthalt kam das \VUgras{Geschwader} und ankerte hinter der
riesigen \VUgras{Mole} \textsuperscript{(4)}, die das \VUgras{Kriegshafenbecken}
\textsuperscript{(5)} abschließt und schützt, und ich konnte in aller Ruhe die
verschiedenen \VUgras{Kriegsschiffe} bewundern, vom kleinen \VUgras{Unterseeboot}
\textsuperscript{(10)}, das taucht und unter Wasser verschwindet, bis zum riesigen
\VUgras{Schlachtschiff} \textsuperscript{(9)}, das bis jetzt kaum etwas zu fürchten
hatte außer den Angriffen des \VUgras{Torpedoboots} \textsuperscript{(8)}.

%%K t10-tit-2 sub
\VUsaut{3.0mm}
\VUpk{10.2pt}{40}{Die kleinen Fahrzeuge.}
\VUsaut{3.0mm}

%%K t10-04-1 p
4. --- Es fand die schwache Stelle im dicken Panzer schon sehr \VUgras{geschickt} und
setzte dort den gefährlichen \VUgras{Torpedo} an, dessen Explosion das Schiff
versenkt; aber es war doch weniger zu fürchten als das Unterseeboot, weil die
Zerstörer es leicht jagen können.

%%K t10-05-1 p
5. --- Ich konnte auch die schnellen \VUgras{Panzerkreuzer} \textsuperscript{(6)}
sehen, fast so unverwundbar wie ein richtiges Schlachtschiff, mehrere
\VUgras{Truppentransporter} \textsuperscript{(12)}, drei oder vier
\VUgras{Kanonenboote} \textsuperscript{(7)} und zuletzt eine \VUgras{Fregatte}
\textsuperscript{(11)}. \VUgras{Am schönsten von allem ist der Anblick dieser Flotte,
wenn sie die Anker lichtet.}

%%K t10-06-1 p
6. --- Was für ein Unterschied im Bau zwischen diesen großen Stahlmassen und den
eleganten \VUgras{Dreimastern} oder den zierlichen \VUgras{Segelschiffen}
\textsuperscript{(13)}, die in der Handelsflotte noch fahren und die ich unter vollen
Segeln in den Hafen einlaufen sah, wenn ich auf der \VUgras{Mole}
\textsuperscript{(14)} spazieren ging. Freilich verloren sie viel von ihrem Vorteil,
wenn der \VUgras{Wind} gegen sie stand. Sie mussten alle Segel bergen, um ins Becken
zu kommen, und ein \VUgras{Schleppdampfer} \textsuperscript{(16)} musste sie holen
und wie schlichte \VUgras{Kähne} \textsuperscript{(17)} an den Kai bringen.

%%K t10-tit-3 sub
\VUsaut{3.0mm}
\VUpk{10.2pt}{40}{Der Handelshafen.}
\VUsaut{3.0mm}

%%K t10-07-1 p
7. --- Was den Hafen, von dem ich spreche, so interessant macht, ist, dass er nicht
nur ein Kriegshafenbecken hat, das durch gewaltige Arbeiten künstlich geschaffen
wurde, sondern auch einen großartigen \VUgras{Handelshafen} an der \VUgras{Mündung}
eines großen Flusses.

%%K t10-08-1 p
8. --- Die Landzunge zwischen den beiden Becken ist befestigt und von einer Reihe
von \VUgras{Batterien} und \VUgras{Bastionen} gedeckt, die ins Meer hinausreichen.
Um auf den \VUgras{Wällen} \textsuperscript{(15)} zu gehen, braucht man eine
besondere Erlaubnis. Ich habe leicht eine bekommen, durch meinen Onkel, der
Ingenieur auf der \VUgras{Werft} \textsuperscript{(18)} ist, und bin die Schiffe
ansehen gegangen, die unter riesigen Hallen gebaut werden.

%%K t10-09-1 p
9. --- Ich habe sogar den Stapellauf eines Schlachtschiffs gesehen, und ich erinnere
mich an das Schaudern, das durch die ganze Menge ging, als die ungeheure Masse, sich
selbst überlassen, ihre Helling hinunterzugleiten begann und auf dem Wasser des
Flusses aufschwamm.

%%K t10-10-1 p
10. --- Zur Werft kommt man über den \VUgras{Exerzierplatz} \textsuperscript{(19)},
wo die Truppen der Garnison jeden Tag üben, und über eine eiserne
\VUgras{Fußgängerbrücke} \textsuperscript{(20)}, die den Exerzierplatz mit dem
\VUgras{Arsenal} \textsuperscript{(21)} verbindet.

%%K t10-tit-4 sub
\VUsaut{3.0mm}
\VUpk{10.2pt}{40}{Das Arsenal und die Docks.}
\VUsaut{3.0mm}

%%K t10-11-1 p
11. --- Mit meinem Onkel habe ich diese große Anlage besichtigt, voller
\VUgras{Geschütze} \textsuperscript{(22)}, \VUgras{Kanonenkugeln}
\textsuperscript{(23)} und \VUgras{Granaten}. Ich habe auch die
\VUgras{Eisenhütte} \textsuperscript{(24)} gesehen, wo die \VUgras{Kessel}
\textsuperscript{(25)} der Dampfer gebaut werden.

%%K t10-12-1 p
12. --- Ganz in der Nähe sind die \VUgras{Docks} \textsuperscript{(26)} --- Trockendocks
--- und die sehr bemerkenswerten \VUgras{Schwimmdocks} \textsuperscript{(27)}, in
denen ich an einem Schiff in Reparatur die \VUgras{Schraube}
\textsuperscript{(28)}, den \VUgras{Kiel} \textsuperscript{(29)} und das
\VUgras{Ruder} \textsuperscript{(30)} eines Fahrzeugs aus der Nähe sehen konnte.

%%K t10-13-1 p
13. --- Der Handelshafen lohnt das Studium ebenso gut wie der Kriegshafen, und ich
habe viele Stunden auf den Kais herumgetrödelt, wo die \VUgras{Raddampfer}
\textsuperscript{(31)} liegen, die den Fluss hinauffahren, und dem
\VUgras{Bockkran} \textsuperscript{(32)} bei der Arbeit zugesehen, der riesige
Lasten hebt, als wären es Federn.

%%K t10-tit-5 sub
\VUsaut{4.0mm}
\VUpk{10.2pt}{40}{Der Lotse.}
\VUsaut{3.0mm}

%%K t10-14-1 p
14. --- Ich habe die \VUgras{Ozeandampfer} \textsuperscript{(34)} bewundert, die mit
wehenden \VUgras{Flaggen} \textsuperscript{(35)} die Reede verließen, gesteuert von
einem geschickten \VUgras{Lotsen} \textsuperscript{(36)}, der wunderbar der
\VUgras{Fahrrinne} folgte, die durch \VUgras{Bojen} \textsuperscript{(40)} und
\VUgras{Baken} bezeichnet ist, und dabei den \VUgras{Leichtern}
\textsuperscript{(41)} auswich, die sich mit ihrem einzigen Rahsegel so schwer
steuern lassen. Am \VUgras{Bug} \textsuperscript{(37)} konnte ich die
\VUgras{Kabinen} \textsuperscript{(39)} der zweiten Klasse erkennen und am
\VUgras{Heck} \textsuperscript{(38)} die Salons der ersten.

%%K t10-15-1 p
15. --- Manchmal habe ich auch zugesehen, wie ein \VUgras{Frachtdampfer}
\textsuperscript{(42)} beladen wurde, in den die Waren aus dem \VUgras{Lagerhaus}
\textsuperscript{(43)} und dem \VUgras{Hafenbahnhof} \textsuperscript{(44)} gestapelt
wurden, die die vielen Züge der \VUgras{Kaibahn} \textsuperscript{(45)}
heranbrachten.

%%K t10-tit-6 sub
\VUsaut{3.0mm}
\VUpk{10.2pt}{40}{Rudern zum Vergnügen.}
\VUsaut{3.0mm}

%%K t10-16-1 p
16. --- Ich bin oft im \VUgras{Hafenbecken} \textsuperscript{(53)} gerudert, wo die
kleinen \VUgras{Boote} \textsuperscript{(46)} Schutz finden, und ich habe das
\VUgras{Ruder} \textsuperscript{(47)} sehr gern geführt. Ich stieg auf das
\VUgras{Deck} \textsuperscript{(48)} eines \VUgras{Segelboots}
\textsuperscript{(49)}, kletterte an der \VUgras{Takelage} \textsuperscript{(52)} den
\VUgras{Mast} \textsuperscript{(51)} hinauf in die \VUgras{Rahen}
\textsuperscript{(50)} und setzte meinen Ausflug dann in einem \VUgras{Skiff}
\textsuperscript{(54)} zwischen den schweren \VUgras{Fischerbooten}
\textsuperscript{(55)} fort, wo die \VUgras{Netze} \textsuperscript{(56)} zum Trocknen
hängen.

%%K t10-17-1 p
17. --- Auf dem \VUgras{Kai} \textsuperscript{(58)} zogen die verschiedensten
\VUgras{Waren} \textsuperscript{(59)} an mir vorbei, in \VUgras{Kisten}
\textsuperscript{(60)} oder in \VUgras{Ballen} \textsuperscript{(61)} verpackt. Sie
bildeten die \VUgras{Ladung} \textsuperscript{(63)} der Kähne, und kräftige
\VUgras{Kräne} \textsuperscript{(62)} hoben sie heraus und setzten sie auf
\VUgras{Lastwagen} \textsuperscript{(64)}, während die \VUgras{Spediteure} sie am
\VUgras{Zollamt} \textsuperscript{(65)} anmeldeten und den Zoll bezahlten.

%%K t10-18-1 p
18. --- Wenn ich schließlich an die \VUgras{Drehbrücke} \textsuperscript{(66)} oder
an die \VUgras{Schleuse} \textsuperscript{(69)} kam, vorbei am \VUgras{Bagger}
\textsuperscript{(68)}, der den \VUgras{Kanal} \textsuperscript{(67)} ausräumt, ging
ich neben dem \VUgras{Signalmast} \textsuperscript{(70)} an der Mole an Land.

%%K t10-19-1 p
19. --- Auf der anderen Seite dieser Mole legen die \VUgras{Barkassen}
\textsuperscript{(71)} an, die die Offiziere des Geschwaders an Land bringen. Es ist
ein schöner Anblick, wenn die Matrosen im Takt die Riemen aufholen, während das Boot
von der \VUgras{Woge} \textsuperscript{(72)} getragen leicht an die Landungstreppe
kommt. Das ist der beste Platz, um die \VUgras{Gezeiten} und den Gang von
\VUgras{Ebbe} und \VUgras{Flut} am \VUgras{Strand} \textsuperscript{(73)} zu
beobachten.

%%K t10-tit-7 sub
\VUsaut{3.0mm}
\VUpk{10.2pt}{40}{Das Offizierskasino.}
\VUsaut{3.0mm}

%%K t10-20-1 p
20. --- Um nach Hause zu kommen, nahm ich die \VUgras{elektrische Straßenbahn}
\textsuperscript{(74)}, deren \VUgras{Haltestelle} \textsuperscript{(75)} am
\VUgras{Pfahl} vor dem Offizierskasino ist.

%%K t10-20-2 p
Vom \VUgras{Balkon} \textsuperscript{(76)} des Kasinos, der mit \VUgras{Ankern}
\textsuperscript{(77)}, Geschützen und Kanonenkugeln geschmückt ist, habe ich oft
eine prächtige Versammlung von Marineoffizieren gesehen: \VUgras{Leutnants}
\textsuperscript{(78)} mit ihren Degen, \VUgras{Artilleriehauptleute}
\textsuperscript{(79)} und \VUgras{Infanteriehauptleute} \textsuperscript{(80)} mit
ihren \VUgras{Säbeln}. Hinter ihnen stand ein \VUgras{Matrose}
\textsuperscript{(81)}, der ihnen ein \VUgras{Fernrohr} brachte, wenn sie sehen
wollten, was in der Ferne vorging.

%%K t10-21-1 p
21. --- Ich habe auch junge \VUgras{Leutnants zur See} \textsuperscript{(82)} gesehen,
die ihre Befehle von ihrem \VUgras{Kommandanten} \textsuperscript{(83)} oder vom
\VUgras{Admiral} \textsuperscript{(84)} entgegennahmen, während ein
\VUgras{Quartiermeister} \textsuperscript{(85)} wartete, um sie an Bord zu tragen.

%%K t10-22-1 p
22. --- Von diesem Balkon ist die Aussicht großartig: er beherrscht den ganzen
Strand mit seinen \VUgras{Zelten} \textsuperscript{(86)} und \VUgras{Badekabinen}
\textsuperscript{(87)}, und zur Badezeit sieht man die \VUgras{Badenden}
\textsuperscript{(88)} vor dem \VUgras{Kasino} \textsuperscript{(89)} fröhlich
herumtollen.

%%K t10-23-1 p
23. --- Am Ende des Strands bildet die Küste eine kleine felsige \VUgras{Halbinsel}
\textsuperscript{(91)}, an der die \VUgras{Brandung} \textsuperscript{(92)} bricht und
auf der ein \VUgras{Leuchtturm} \textsuperscript{(93)} steht, der den Seeleuten
nachts den Weg weist. Diese Halbinsel hängt nur durch eine sehr schmale
\VUgras{Landenge} \textsuperscript{(90)} mit dem Land zusammen.

%%K t10-tit-8 sub
\VUsaut{3.0mm}
\VUpk{10.2pt}{40}{Ein Blick über die Küste.}
\VUsaut{3.0mm}

%%K t10-24-1 p
24. --- Weiter zieht sich die \VUgras{Steilküste} \textsuperscript{(94)} hin, vom Meer
angenagt, das ihr eine launische Folge von \VUgras{Buchten} \textsuperscript{(95)} und
\VUgras{Landspitzen} (Kaps) ausschneidet und sie damit sehr malerisch macht.

%%K t10-24-2 p
Auf der \VUgras{Hochfläche} \textsuperscript{(96)} über der \VUgras{Meerenge}
\textsuperscript{(98)} stehen ein sehr wichtiges \VUgras{Fort} \textsuperscript{(97)}
und ein paar Villen.

%%K t10-25-1 p
25. --- Diese Meerenge trennt vom Festland eine sehr gefährliche \VUgras{Insel}
\textsuperscript{(99)}, von \VUgras{Felsen} \textsuperscript{(100)} und
\VUgras{Riffen} umgeben, an denen manchmal Boote und sogar große Schiffe scheitern.
So schnell die Rettung auch kommt und so rasch die \VUgras{Rettungsboote} auch
ausfahren, diese \VUgras{Schiffbrüche} sind schrecklich, und in den
\VUgras{Äquinoktialstürmen} sind mehrere Schiffe mit Mann und Ladung untergegangen.

%%K t10-25-2 p
Trotz dieser Nachbarschaft wird \VUgras{der Hafen} von den Seeleuten viel benutzt.
\VUsaut{6.0mm}
\VUfilet{22mm}
